%! Simplifying Algebraic Expressions and Exponent Rules — Unit 1, Lesson 1.1 — Lesson Plan
\documentclass[10pt]{article}
\usepackage{atda-article}
\usepackage{atda-boxes}
\usepackage{graphicx}   % -article does NOT load graphicx; needed for thumbnails/screenshots
\graphicspath{{images/}}
\newcommand{\TallMath}[1]{$\displaystyle #1\rule[-1.4em]{0pt}{3.2em}$}
\newcommand{\UnitNumberName}{Unit 1: Algebraic Expressions and Factoring \quad}
\newcommand{\LessonNumberName}{Lesson 1.1: Simplifying Algebraic Expressions and Exponent Rules}

\begin{document}
\begin{center}
    {\Large\bfseries \CourseName \\
    {\small\bfseries \UnitNumberName \LessonNumberName}}
\end{center}

\begin{tcolorbox}[colback=mist, colframe=royal, boxrule=0.9pt, arc=2mm,
    left=2mm, right=2mm, top=1.5mm, bottom=1.5mm]
    \textbf{Primary Objective:} I can simplify monomial expressions in one and two variables with
    the exponent rules, add and subtract polynomials by combining like terms, and multiply
    polynomials so that every term of one factor reaches every term of the other. I can build a
    polynomial from a situation --- revenue, cost, profit, volume --- evaluate it, and say what
    the number means.
\end{tcolorbox}

\renewcommand{\arraystretch}{1.6}
\begin{skillbox}[Priority Ideas \& Skills]{goldbox}
\begin{tabularx}{\linewidth}{@{}X|X@{}}
\begin{minipage}[t]{\linewidth}\vspace{0pt}
\textbf{Knowledge \& Skills covered --- \texttt{A2.EO.3a}:}
\begin{itemize}[leftmargin=1.1em, itemsep=2pt, topsep=2pt]
  \item Determine \textbf{sums} and \textbf{differences} of polynomials in one and two variables,
        writing the result in standard form.
  \item Determine \textbf{products} of polynomials: monomial $\times$ polynomial, binomial
        $\times$ binomial (one and two variables), binomial $\times$ trinomial.
\end{itemize}
\textbf{Supporting fluency (Algebra 1--2 review, not a new code):}
\begin{itemize}[leftmargin=1.1em, itemsep=2pt, topsep=2pt]
  \item The exponent rules: product, quotient, power of a power, power of a product, power of a
        quotient, and the zero and negative exponents.
  \item Combining like terms; degree and standard form in one and two variables.
\end{itemize}
\textbf{Where these codes go next:} \texttt{A2.EO.2a--c} in Lesson 1.2 (the same rules extended
to fractional exponents); \texttt{A2.EO.3b} in 1.3--1.5, which is this lesson's multiplication
run backwards; \texttt{A2.F.1} in Unit 3, where the product rule reappears as
$\log(ab)=\log a+\log b$.
\end{minipage}
&
\begin{minipage}[t]{\linewidth}\vspace{0pt}
\textbf{Key Understandings (the \emph{why}):}
\begin{itemize}[leftmargin=1.1em, itemsep=2pt, topsep=2pt]
  \item \textbf{An exponent counts factors --- every rule is bookkeeping on that count.} Written
        out, $x^3\cdot x^4$ is seven $x$'s. Students who can reconstruct a rule in five seconds
        never need to guess between adding and multiplying exponents.
  \item \textbf{Coefficients multiply; exponents add.} These are two different operations happening
        in one expression, and confusing them is the single most common error of the lesson.
  \item \textbf{Zero and negative exponents are forced, not chosen.} $x^0=1$ and $x^{-n}=1/x^n$ are
        what the quotient rule \emph{requires}; presenting them as definitions to memorize is what
        makes students believe $x^{-2}$ is negative.
  \item \textbf{Subtraction is distribution in disguise.} $-(a+b)$ is $-a-b$. The minus sign in
        front of a polynomial reaches every term, exactly as the $-16$ did in Lesson 1.0.
  \item \textbf{Products are exhaustive, not patterned.} FOIL is a special case, not a rule ---
        ``every term times every term'' scales to a trinomial and FOIL does not.
  \item \textbf{Polynomials come from situations.} Revenue $-$ cost is a subtraction students can
        check against the story; that check is what makes the algebra worth doing.
\end{itemize}
\end{minipage}
\end{tabularx}
\end{skillbox}

\begin{skillbox}[{Vocabulary, Concepts, \& Theorems}]{greenbox}
\renewcommand{\arraystretch}{1.35}
\begin{tabularx}{\linewidth}{@{}p{3.6cm}X@{}}
\textbf{Term} & \textbf{Meaning students record} \\
\hline
Base and exponent & In $x^7$, $x$ is the base and $7$ is the exponent. The exponent counts the
    factors of the base being multiplied. \\
Like terms & Terms with exactly the same variables raised to exactly the same powers. $5m^2$ and
    $-2m^2$ combine; $5m^2$ and $7m$ do not. \\
Degree & The largest total exponent in any one term. In two variables, add the exponents within
    a term: $3a^2b^3$ has degree $5$. \\
Standard form & Terms ordered from highest degree to lowest, ending with the constant. \\
Distributive property & $a(b+c)=ab+ac$. The reason a factor outside parentheses reaches every
    term inside --- and the engine behind every product in this lesson. \\
Monomial / binomial / trinomial & Polynomials with one, two, and three terms. Naming them makes
    ``a binomial times a trinomial gives six products'' a usable check. \\
Product of powers & $a^m\cdot a^n=a^{m+n}$. Quotient: $a^m/a^n=a^{m-n}$. Power of a power:
    $(a^m)^n=a^{mn}$. Power of a product: $(ab)^n=a^nb^n$. \\
Zero \& negative exponent & $a^0=1$ and $a^{-n}=1/a^n$, for $a \ne 0$. Both follow from the
    quotient rule; neither makes a value negative. \\
\end{tabularx}
\end{skillbox}

\begin{fixedskillbox}[Activate Prior Knowledge \& Spiral Review]{mist}
\begin{tabularx}{\linewidth}{@{}X c@{}}
\begin{minipage}[t]{0.55\linewidth}\vspace{0pt}
\textbf{Four items, about 6 minutes:}
\begin{enumerate}[leftmargin=1.4em, itemsep=1pt, topsep=2pt]
  \item Write $x^3\cdot x^4$ out as factors and count them (\emph{discovers} the product rule)
  \item Combine like terms in $5m^2-3m+8-2m^2+7m$
  \item Distribute and simplify $-3(2y-5)+4y$
  \item Spiral from 1.0: expand $A(x)=x(30-3x)$
\end{enumerate}
\end{minipage}
&
\begin{minipage}[t]{0.38\linewidth}\vspace{0pt}
\textbf{How to use the results:} item 1 is not a review question --- it is the lesson's first two
minutes. Take answers to it \emph{before} the hook and write the seven $x$'s on the board; the
product rule should be stated by a student, not by you. Items 2--4 are diagnostic for anyone who
missed warm-up items 1--2 on the Unit 1 opener.
\end{minipage}
\end{tabularx}
\end{fixedskillbox}

\begin{skillbox}[Hook]{mist}
\textbf{The Party Size tub.} A theater's regular popcorn tub is a cylinder $3$ inches in radius
and $8$ inches tall. The \emph{Party Size} is twice as wide and twice as tall --- radius $6$,
height $16$ --- and costs \textbf{three times} as much. \textbf{Pose it this way:} ``Twice as wide,
twice as tall, three times the price. Good deal or bad deal? Vote before you compute.'' Most rooms
vote ``bad deal,'' reasoning that twice the size should be about twice the popcorn. Then compute
with $V=\pi r^2h$:
\[ \pi(2r)^2(2h) = \pi\cdot 4r^2\cdot 2h = 8\pi r^2h. \]
Eight times the popcorn for three times the price. Concretely: $72\pi \approx 226$ cubic inches
becomes $576\pi \approx 1810$. The exponent on $r$ is what the intuition missed --- doubling a
squared quantity quadruples it. \textbf{That is the whole lesson in one line:} the exponent rules
are not notation, they are how quantities actually scale.
\end{skillbox}

\boxguard
\begin{skillbox}[Lesson]{mist}
\begin{multicols}{2}
\textbf{1. Build the rules by counting (10 min).} Start from the warm-up's seven $x$'s and fill
the notes table together. Do \emph{not} present five rules to memorize --- derive each one by
writing factors out, then let students state it. \textbf{Ask:} ``Why does $(2d)^4$ start with
$16$?'' The exponent is outside the parentheses, so it reaches the $2$ as well. Work the two
examples, naming the coefficients-multiply/exponents-add split out loud.

\textbf{2. Force the zero and negative exponents (8 min).} $\dfrac{x^5}{x^5}$ is $x^0$ by the
quotient rule and $1$ by cancelling, so $x^0=1$. $\dfrac{x^3}{x^7}$ is $x^{-4}$ and $\dfrac{1}{x^4}$,
so those are equal. \textbf{Ask:} ``Is $2^{-3}$ negative?'' Get the answer $\tfrac18$ and the
sentence ``a negative exponent moves the factor across the fraction bar.'' This is the
misconception to kill today; it costs points in Lesson 1.2 and again in Unit 3.

\textbf{3. Sums and differences (7 min).} Do the two generic examples side by side so the only
difference visible is the sign handling. \textbf{Ask:} ``What happened to the $-9$?'' Then build
the concession-stand model: revenue, cost, profit. Have students check $P(8)=380$ against the
story ($70$ trays $\times$ \$8 $=$ \$560, minus \$180 in costs) --- that verification is the habit
worth more than the answer.

\textbf{4. Products (10 min).} State the one rule --- every term times every term --- and use the
$2\times 3$ product count as a check. Work all three examples; the two-variable one is where to
slow down, because $-15ab$ and $+2ab$ combining surprises students who think $ab$ terms cannot
combine at all. Name FOIL as a special case that does not extend, then move on.

\textbf{5. Guided practice (8 min).} The tray problem. Multiply the two binomials first, then
distribute the $x$; students who start with $x(10-2x)$ get there too but with messier bookkeeping.
Item 3 (why $x<4$) is the interpretation item --- do not let it get skipped for time.
\end{multicols}
\end{skillbox}

\boxguard
\begin{skillbox}[Explicit Instruction: Multiplying Any Two Polynomials]{mist}
\begin{tabularx}{\linewidth}{@{}X|X@{}}
\begin{minipage}[t]{\linewidth}\vspace{0pt}
\textbf{The four steps students record:}
\begin{enumerate}[leftmargin=1.4em, itemsep=2pt, topsep=2pt]
  \item Count the products you expect: (terms in the first) $\times$ (terms in the second).
  \item Multiply every term of the first by every term of the second. Coefficients multiply;
        exponents add.
  \item Check that you wrote down that many products before combining anything.
  \item Combine like terms and write the result in standard form.
\end{enumerate}
\textbf{Say this out loud:} ``FOIL is what this rule looks like when both factors happen to have
two terms. It is not the rule.''
\end{minipage}
&
\begin{minipage}[t]{\linewidth}\vspace{0pt}
\textbf{Worked example (two variables):}
\[
\begin{aligned}
(3a+2b)(a-5b) &= 3a^2-15ab+2ab-10b^2 \\
              &= 3a^2-13ab-10b^2
\end{aligned}
\]
Four products, then two of them combine. \textbf{The error to name before it happens:} treating
$-15ab$ and $+2ab$ as unlike terms because two letters are involved. They have the same variables
to the same powers, so they combine --- the answer has three terms, not four.
\end{minipage}
\end{tabularx}
\end{skillbox}

\begin{skillbox}[Active Monitoring]{redbox}
\begin{itemize}[leftmargin=1.2em, itemsep=2pt, topsep=2pt]
  \item \textbf{Notes, box 1:} watch for $(2m^3n)^4$ answered as $2m^{12}n^4$ --- the coefficient
        left untouched. Ask ``what is the base of that fourth power?''
  \item \textbf{Notes, box 2:} watch for $4a^4b^{-3}$ written as $-\dfrac{4a^4}{b^3}$. The sign of
        the exponent is not the sign of the value.
  \item \textbf{Notes, box 3:} the giveaway error is $-10x^2+150x-20x+340$ --- the minus reached
        only the first term. Point at the parentheses, do not supply the fix.
  \item \textbf{Activity, Tier A item 3:} the second polynomial is $-9x+195$, so subtracting it
        \emph{adds} $9x$. Expect this to be the sticking point for a full third of the room.
  \item \textbf{Cold-call prompts:} ``Why is the Party Size eight times, not four?'' ``Is $3^{-2}$
        negative?'' ``How many products before you combine, and how do you know?''
\end{itemize}
\end{skillbox}

\boxguard
\begin{skillbox}[Group Work \& Differentiation]{redbox}
\begin{multicols}{3}
\textbf{Tier R --- Remediate}
\begin{itemize}[leftmargin=1em, itemsep=1pt, topsep=2pt]
  \item Five one-line fluency items: a monomial product, a power of a product, a quotient, a
        polynomial sum, and a monomial times a trinomial.
  \item Support: have the group write the factors out longhand on the first one.
\end{itemize}

\columnbreak
\textbf{Tier A --- Approaching Proficiency}
\begin{itemize}[leftmargin=1em, itemsep=1pt, topsep=2pt]
  \item The school-store model: build $R(x)$, $C(x)$, then $P(x)=R-C$ --- the minus-sign item.
  \item Evaluate $P(6)=183$ and interpret it in context.
  \item A two-variable product, $(2x+5y)(3x-y)$.
\end{itemize}

\columnbreak
\textbf{Tier E --- Extension}
\begin{itemize}[leftmargin=1em, itemsep=1pt, topsep=2pt]
  \item Critique two wrong answers, $(4x^3)^2$ and $(2m+3)^2$, and name the rule each breaks.
  \item Scaling: a cube of edge $3s$ has $27$ times the volume but only $9$ times the surface.
  \item Justify why $P(0)=-195$ is correct, not an arithmetic error.
\end{itemize}
\end{multicols}
\end{skillbox}

\begin{skillbox}[Individual Work \& Assessment]{redbox}
\textbf{Exit ticket (3 items, no notes):} (1) simplify $\dfrac{12a^5b^3}{3a^2b^7}$ with no negative
exponents --- answer $\dfrac{4a^3}{b^4}$; (2) subtract $\left(5x^2-3x+7\right)-\left(2x^2+4x-9\right)$
--- answer $3x^2-7x+16$; (3) \textbf{the conceptual check} --- a classmate writes
$\left(3x^2\right)^3=3x^6$; give the correct answer and say in one sentence what was forgotten.

\textbf{Using the results:} sort into three piles --- \emph{clean}, \emph{one rule slipping},
\emph{cannot start}. Item 1 isolates the negative exponent, item 2 isolates the distributed minus,
item 3 isolates the coefficient under a power. Because each item tests a different failure, the
pile a student lands in names the fix, and that is your Tier R grouping for Lesson 1.2.
\end{skillbox}

\boxguard[18]
\begin{skillbox}[Reinforcement \& Extension]{goldbox}
\textbf{Homework} (12 items): seven fluency items covering all three exponent-rule families and
all three polynomial operations, a three-part concession-stand profit model paralleling the
activity with new numbers, and two look-ahead items --- a cube-scaling justification and a square
whose area is $49x^6$.

\textbf{Extension (optional):} simplify the surface-area-to-volume ratio of a cube, $\dfrac{6s^2}{s^3}
= \dfrac{6}{s}$, and use it to explain why one large block of ice outlasts the same ice in small
cubes. It is the hook's scaling idea turned into a quotient.

\textbf{Preview of Lesson 1.2 (\texttt{A2.EO.2a--c}):} radicals and rational exponents. Homework
item 12 asks students to undo a square; Lesson 1.2 gives that move a symbol and shows the same
exponent rules survive when the exponents become fractions.

\textbf{Connections.} Covers \texttt{A2.EO.3a} (sums, differences, and products of polynomials in
one and two variables), supported by Algebra 1--2 exponent-rule review. Builds on Lesson 1.0's
expand-and-compare; feeds \texttt{A2.EO.2a--c} (1.2), \texttt{A2.EO.3b} (1.3--1.5, this lesson
reversed), and \texttt{A2.F.1} in Unit 3, where the product rule becomes a property of logarithms.
\end{skillbox}

% ── Teacher notes — one per component, in packet order ────────────────────────
% Teacher-only prose lives HERE, never in a _key: a note is the one block in a
% key with no counterpart in the blank, so it makes the key longer than the
% blank. See references/conventions.md ("teachernote").
\begin{teachernote}[Warm-Up]
    \textbf{6 minutes, then take item 1 out loud before anything else.} Answers: (1) $x^7$, after
    writing seven $x$'s; (2) $3m^2+4m+8$; (3) $-2y+15$; (4) $-3x^2+30x$. Item 1 is the pivot of
    the whole lesson --- if you score it silently and move on, you have spent the discovery and
    bought nothing. Put the seven $x$'s on the board and ask who can state the rule. Items 2 and 3
    tell you who is still losing signs; item 4 is the Lesson 1.0 spiral, and a student who cannot
    expand $x(30-3x)$ tonight will not survive the profit model in the activity.
\end{teachernote}

\begin{teachernote}[Guided Notes]
    \textbf{The table in box 1 is a fill-in, not a handout --- resist reading the five rules
    aloud.} Answers across the ``Try it'' column: $x^{13}$, $y^7$, $c^{15}$, $16d^4$,
    $\dfrac{k^2}{9}$. The three blanks elsewhere are \emph{coefficient}, \emph{negative}, and
    \emph{every}; box 4's count is \emph{six}. Box 2 is short but it is the box that pays off in
    Lesson 1.2 --- do not cut it for time. If the period is running behind, cut the second worked
    product in box 4 (the two-variable one is the keeper) rather than the tray in guided practice.
    On the tray, expect at least one group to multiply $x(10-2x)$ first; that is not wrong, so let
    them finish and compare the two routes at the board.
\end{teachernote}

\begin{teachernote}[Group Activity]
    \textbf{Groups of three, 15 minutes, all groups start at Tier R.} Tier R is five one-liners;
    a group still working on it at six minutes needs you, not more time. Tier A item 3 is the
    lesson's real assessment: subtracting $-9x+195$ requires adding $9x$, and roughly a third of
    the room will write $-9x$ instead. Do not correct it --- ask them to check $P(6)$ against the
    story ($54$ bottles at \$6 is \$324; costs are \$60 $+$ \$81 $=$ \$141; profit \$183) and let
    the mismatch send them back. Tier E item 3 is the one to share out: a negative profit at a
    price of zero is not a broken model, it is a stand that gave away $90$ bottles and still owes
    \$195.
\end{teachernote}

\begin{teachernote}[Exit Ticket]
    \textbf{5 minutes, notes closed, hand it to me at the door.} Each item isolates one failure
    mode on purpose --- negative exponent, distributed minus, coefficient under a power --- so
    read the tickets by \emph{which} item is wrong, not by a score out of three. A student who
    misses only item 1 needs five minutes on notation; a student who misses only item 2 is losing
    signs and will lose them again in 1.3 when factoring out a negative GCF. Item 3 asks for a
    sentence: accept anything that says the exponent applies to the $3$, and do not require the
    phrase ``power of a product.''
\end{teachernote}

\begin{teachernote}[Homework]
    \textbf{Expect 20--25 minutes.} Items 1--7 are deliberately unmixed by type so a student can
    tell you which \emph{kind} of problem broke, not just that homework was hard. Item 3 is the
    only one with both a negative and a zero exponent, so check it first when you spot-grade.
    Items 8--10 repeat the activity's model with new numbers, on purpose: a group that leaned on
    its partners in class has to carry it alone the same night. Item 12 is the bridge into Lesson
    1.2 --- students are asked to \emph{check} $7x^3$ by squaring rather than to find a square
    root, so nobody needs a symbol they have not met. The extension is genuinely optional; it is
    the best follow-up for whoever finished Tier E early.
\end{teachernote}
\end{document}
